\chapter{Applications et Active Directory}
\label{ch:apps-ad}

Ce chapitre détaille ce qui tourne dans le lab côté applicatif et
identités. C'est là qu'il y a de quoi auditer~: applications avec
des vulnérabilités plantées exprès, AD avec des utilisateurs et
des groupes représentatifs.

\section{CorpNet : le portail interne}

\texttt{CorpNet} est le premier service applicatif déployé sur
\texttt{vm-web01}. C'est un portail web en \textbf{PHP} très
classique~: page de login, mur d'actualités, agenda, chat
flottant, espace personnel.

Il n'est pas \emph{secure by design}. Il a été développé avec des
vulnérabilités intentionnelles pour servir de cible aux scénarios
d'attaque et donner à Wazuh des choses à détecter~:

\begin{itemize}
  \item \textbf{XSS} stockée dans le chat flottant~: un message
    JavaScript posté est exécuté chez tous les destinataires.
  \item \textbf{IDOR} sur le polling du chat~:
    \texttt{?conversation\_id=N+1} renvoie les messages d'autres
    utilisateurs.
  \item \textbf{Injection SQL} sur certaines pages de recherche
    (jointures non paramétrées).
  \item \textbf{Mots de passe stockés en clair} dans la table
    \texttt{users}.
\end{itemize}

\section{MaFormation et MaCandidature : applications métier}

Pour aller au-delà d'un seul service, deux applications métier
ont été ajoutées dans la stack \texttt{internal-apps}~:

\begin{description}
  \item[MaFormation.] Plateforme pédagogique~: cours, inscriptions,
    notes. Codée en Node.js / Express, base PostgreSQL, ORM
    Prisma. C'est l'application \emph{secure by design}~:
    container non-root, filesystem en read-only, capacités
    Linux droppées, communication avec la base sur un network
    Docker \texttt{internal: true}.
  \item[MaCandidature.] Plateforme de candidature aux formations~:
    création de compte, dépôt de dossier, suivi. Même stack
    technique que MaFormation, \textbf{mais une vulnérabilité
    intentionnelle~:} le socket Docker \texttt{/var/run/docker.sock}
    est monté dans le container, et le container tourne en root
    avec le CLI \texttt{docker} installé.
\end{description}

Cette vulnérabilité côté MaCandidature est précisément le genre
de chose qu'un dossier d'homologation est censé documenter et
faire corriger. C'est le scénario \emph{«~ce que mon
dossier devrait dire, mais ne dit plus si je ne le mets pas à
jour~»}.

\paragraph{Référence.} CIS Docker Benchmark 5.31~; MITRE
ATT\&CK T1611 (\emph{Escape to Host}).

\section{Architecture Docker sur vm-web01}

\begin{figure}[h]
\centering
\begin{small}
\begin{verbatim}
vm-web01 (snet-dmz) -- Docker engine
  +- stack "corpnet"           : portail PHP + MySQL
  +- stack "internal-apps"
      +- nginx (TLS, loopback)
      |    +- :8443 -> maformation_app
      |    +- :8444 -> macandidature_app
      +- maformation_app  + maformation_db
      +- macandidature_app (*) docker.sock monte
          + macandidature_db
\end{verbatim}
\end{small}
\caption{Stacks Docker sur le serveur web.}
\label{fig:docker-web01}
\end{figure}

\paragraph{Trois networks Docker isolés.}
\texttt{maformation\_proxy\_net}, \texttt{macandidature\_proxy\_net}
pour le trafic nginx~$\leftrightarrow$~app, et deux networks
\texttt{internal: true} pour les bases PostgreSQL (zéro accès
Internet, même depuis Docker).

\paragraph{nginx en frontal.}
Il termine TLS, applique HSTS, CSP, X-Frame-Options, fait du
rate-limit. Il n'écoute qu'en loopback (\texttt{127.0.0.1:8443}
et \texttt{:8444}) sur vm-web01~; l'exposition vers l'extérieur
passe par CorpNet en reverse-proxy.

\paragraph{Secrets.}
Les mots de passe (base, JWT, bind LDAP) sont dans des fichiers
Docker secrets sous \texttt{secrets/}, jamais en variables
d'environnement clair.

\section{Active Directory : corpnet.local}

Le contrôleur \texttt{vm-dc01} héberge le domaine
\textbf{\texttt{corpnet.local}}, créé via AD~DS sur Windows Server
2022. Le domaine est volontairement réaliste pour un ENT de
moyenne taille~:

\begin{itemize}
  \item Une unité d'organisation (\emph{OU}) par direction~:
    \texttt{OU=Direction}, \texttt{OU=DSI}, \texttt{OU=Pedagogie},
    \texttt{OU=Scolarite}, etc.
  \item Une vingtaine d'utilisateurs (\texttt{m.martin},
    \texttt{p.bernard}, \texttt{j.dupont}, \texttt{s.david}\ldots)
    avec un mot de passe par défaut \texttt{Welcome2025!} qui
    sert à la fois aux tests applicatifs et aux scénarios
    d'attaque (mot de passe faible).
  \item Des groupes de sécurité (\texttt{G\_Etudiants},
    \texttt{G\_Enseignants}, \texttt{G\_Admin}) utilisés pour
    l'autorisation dans les applications.
  \item Un compte de service \texttt{svc\_ldap} pour le bind LDAP
    depuis MaFormation et MaCandidature.
\end{itemize}

\paragraph{Pourquoi un AD réaliste.}
Sans utilisateurs, sans groupes, sans changements, un dossier
d'homologation côté identités n'a rien à dire. L'AD du lab fournit
de la matière~: des comptes inactifs à détecter, des comptes
privilégiés à inventorier, des mots de passe faibles à signaler.

\section{Authentification dans les applications}

Les trois applications (CorpNet, MaFormation, MaCandidature)
authentifient les utilisateurs contre \texttt{corpnet.local} via
LDAP, en utilisant le compte de bind \texttt{svc\_ldap}.

À la première connexion, l'utilisateur est créé localement dans
la base de l'application, avec un rôle déduit du
\texttt{memberOf} LDAP. Cette logique permet de tracer
côté application les attributions de rôles, ce qui est exigé par
l'homologation (qui a accès à quoi, et pourquoi).

\begin{warning}
Toutes les vulnérabilités décrites dans ce chapitre (XSS, IDOR,
SQLi, mot de passe par défaut faible, socket Docker exposé) sont
\textbf{intentionnelles} et confinées au lab. Elles ne doivent
pas être présentes dans un SI en production~; elles sont là
pour que le dossier d'homologation ait quelque chose à dire et
que Wazuh ait quelque chose à détecter.
\end{warning}
